\documentclass[../../main.tex]{subfiles}
\graphicspath{{../../image/}{../../figure/}}
\begin{document}

\section{变差的界}\label{section:ch16-变差界}

遵循\cref{section:ch16-阻尼重分布}中勾勒的程序，我们在此通过建立估计\eqref{align:ch16-s03-e03}来完成\cref{theorem:ch16-1631}的证明。

对固定的初始数据 $U_0$（它属于 $\mathscr{C}_\omega$ 并满足\eqref{align:ch16-s01-e10}，其中 $\omega$ 与 $\delta$ 小），记 $U$ 为\eqref{align:ch16-s02-e01}、\eqref{align:ch16-s01-e02}的（局部）可容 $BV$ 解，定义在某个带形 $(-\infty,\infty)\times[0,T)$ 上。目标是推导函数
\begin{align}
X(t)=TV_{(-\infty,\infty)}U(\cdot,t),\qquad 0\leqslant t<T
\label{align:ch16-s04-e01}
\end{align}
的界。

解 $U$ 通过\eqref{align:ch16-s03-e06}与\eqref{align:ch16-s03-e07}诱导出作为 $(x,t)$ 在 $(-\infty,\infty)\times[0,T)$ 上函数的 $\Phi$ 与 $Z$，从而我们可以把由\eqref{align:ch16-s03-e11}与\eqref{align:ch16-s03-e12}定义的 $\hat F$ 与 $\hat P$ 视为变量 $(\hat U,x,t)$ 的已知函数。计划是把 $\hat U$（重新）构造为非齐次平衡律组\eqref{align:ch16-s03-e10}在由\eqref{align:ch16-s03-e15}定义的初值 $\hat U_0$ 下的解。

$\hat U$ 的构造将借助随机选取法并结合算子分裂来实现。我们的系统具有\eqref{align:ch16-s01-e01}的形式，但它有特殊结构：非齐次性通过函数 $\Phi$ 与 $Z$ 隐式进入。此外，对固定的 $\hat U$，函数 $\hat F$ 与 $\hat P$ 仅仅 Lipschitz。这些特殊特征将迫使我们采用\cref{section:ch16-cauchy问题}中勾勒的一般算法的一个笨重变形。

如同\cref{section:ch16-cauchy问题}那样，我们不失一般性地假设没有任何特征速度为零，固定一个随机序列 $\mathscr{P}=\{a_0,a_1,\ldots\}$（$a_s\in(-1,1)$），取空间步长 $h$ 与时间步长 $\tau=\lambda^{-1}h$，并标识交错网格：网格点 $(x_r,t_s)$，其中 $x_r=rh$、$t_s=s\tau$，$r+s$ 为偶数；采样点 $(y_s^r,t_s)$，其中 $y_s^r=x_r+a_sh$，$r+s$ 为奇数。

约似解 $\hat U_h$ 将在带形
\begin{align}
\mathscr{S}_s=\{(x,t):-\infty<x<\infty,\ t_s\leqslant t<t_{s+1}\},\qquad s=0,1,\ldots,s^*,
\label{align:ch16-s04-e02}
\end{align}
上依次确定，其中 $s^*$ 是使 $(s^*+1)\tau\leqslant T$ 的最大整数。

算法由置 $\hat U_h(x,0)=\hat U_0(x)$（$-\infty<x<\infty$）启动。

现假设 $\hat U_h$ 已在 $\bigcup_{k=0}^{s-1}\mathscr{S}_k$ 上确定，我们通过下述程序把它的定义域延拓到 $\mathscr{S}_s$。

对 $r=0,\pm1,\pm2,\ldots$ 与 $s=0,1,\ldots,s^*$，定义
\begin{align}
\Phi_s^r=\Phi(x_r,t_s),\qquad Z_s^r=Z(x_r,t_s)
\label{align:ch16-s04-e03}
\end{align}
然后，对 $r+s$ 为奇数者，令
\begin{align}
\hat U_s^r=\hat U_h(y_s^r,t_s-),\qquad U_s^r=\hat U_s^r+\Phi_s^r,
\label{align:ch16-s04-e04}\\
Q_s^r=DF(U_s^r)^{-1}DF(\Phi_s^r),
\label{align:ch16-s04-e05}\\
\hat P_s^r=\hat P(\hat U_s^r,\Phi_s^r,Z_s^r),
\label{align:ch16-s04-e06}\\
\begin{cases}
V_s^r=U_s^r+Q_s^r[\Phi_s^{r-1}-\Phi_s^r]-\tau\hat P_s^r,\\
W_s^r=U_s^r+Q_s^r[\Phi_s^{r+1}-\Phi_s^r]-\tau\hat P_s^r,
\end{cases}
\label{align:ch16-s04-e07}\\
\begin{cases}
\hat V_s^r=V_s^r-\Phi_s^{r-1},\\
\hat W_s^r=W_s^r-\Phi_s^{r+1}.
\end{cases}
\label{align:ch16-s04-e08}
\end{align}

接着，对 $r=0,\pm1,\pm2,\ldots$ 与 $s=0,1,\ldots,s^*$（$r+s$ 为偶数者），我们在矩形
\begin{align}
\mathscr{R}_s^r=\{(x,t):x_{r-1}<x<x_{r+1},\ t_s\leqslant t<t_{s+1}\}
\label{align:ch16-s04-e09}
\end{align}
上定义 $\hat U_h$ 为
\begin{align}
\hat U_h(x,t)=\tilde U(x-x_r,t-t_s)-\Phi_s^r,
\label{align:ch16-s04-e10}
\end{align}
其中 $\tilde U$ 是 Riemann 问题
\begin{align}
\partial_t\tilde U+\partial_xF(\tilde U)=0,
\label{align:ch16-s04-e11}\\
\tilde U(x,0)=
\begin{cases}
W_s^{r-1},&-\infty<x<0,\\
V_s^{r+1},&0<x<\infty
\end{cases}
\label{align:ch16-s04-e12}
\end{align}
的解。

为看清上述构造的动机，首先注意：由\eqref{align:ch16-s04-e08}、\eqref{align:ch16-s04-e10}、\eqref{align:ch16-s04-e11}与\eqref{align:ch16-s04-e12}，$\hat U_h$ 在矩形 $\mathscr{R}_s^r$ 上满足方程
\begin{align}
\partial_t\hat U_h+\partial_x\hat F(\hat U_h,\Phi_s^r)=0,
\label{align:ch16-s04-e13}
\end{align}
连同初始条件
\begin{align}
\hat U_h(x,t_s)=
\begin{cases}
\hat W_s^{r-1},&x_{r-1}<x<x_r,\\
\hat V_s^{r+1},&x_r<x<x_{r+1},
\end{cases}
\label{align:ch16-s04-e14}
\end{align}
沿 $\mathscr{R}_s^r$ 的底边。于是，\eqref{align:ch16-s04-e13}是齐次守恒律组，它由\eqref{align:ch16-s03-e10}丢弃源项 $\hat P$ 并把 $\hat F$ 中的 $\Phi$ 冻结在网格点 $(x_r,t_s)$ 处的值 $\Phi_s^r$ 而得到。另一方面，由\eqref{align:ch16-s04-e07}、\eqref{align:ch16-s04-e08}中 $\hat V_s^{r+1}$ 与 $\hat W_s^{r-1}$ 的构造，\eqref{align:ch16-s04-e14}计入了源项的效应以及由 $\Phi$ 与 $Z$ 诱导的非齐次性（算子分裂）。

如同第十三章那样，我们此处的任务是建立算法的紧性与相容性。
紧性将通过一致（关于 $t$ 与 $h$）约似解族 $\{\hat U_h(\cdot,t)\}$ 的总变差的界来展示。我们将在前提 $\omega$ 与 $\hat U_h(\cdot,t)$ 的变差被一个小正常数 $\rho$ 界住之下操作。这将在证明\eqref{align:ch16-s03-e03}之后后验地验证，只要 $\omega$ 与 $\delta$ 充分小。为避免符号激增，我们通篇用 $c$ 表示一个与 $\omega,\delta,\rho,T,h$ 无关的一般正常数。

显然，对任意 $t\in(t_s,t_{s+1})$，$s=0,1,\ldots,s^*$，$\hat U_h(\cdot,t)$ 的总变差取相同值：
\begin{align}
TV_{(-\infty,\infty)}\hat U_h(\cdot,t)=\hat X_s,\qquad t_s<t<t_{s+1}.
\label{align:ch16-s04-e15}
\end{align}
此外，
\begin{align}
\hat X_s=J_s+K_s,\qquad s=0,1,\ldots,s^*,
\label{align:ch16-s04-e16}
\end{align}
其中 $J_s$ 是 $r+s$ 为奇数时由 $\hat U_h$ 越过 $x_r$ 的跳跃生成的那部分变差，$K_s$ 是 $r+s$ 为偶数时由从网格点 $(x_r,t_s)$ 发出的波扇生成的那部分变差。

在 $x=x_r$ 处（$r+s$ 为奇数），$\hat U_h$ 从 $\hat V_s^r$ 跳到 $\hat W_s^r$，跳跃由\eqref{align:ch16-s04-e07}与\eqref{align:ch16-s04-e08}确定：
\begin{align}
\hat W_s^r-\hat V_s^r=[Q_s^r-I][\Phi_s^{r+1}-\Phi_s^{r-1}].
\label{align:ch16-s04-e17}
\end{align}
由于 $U_0\in\mathscr{C}_\omega$，$U$ 满足\eqref{align:ch16-s03-e01}，对 $t\in[0,T)$，由此
\begin{align}
TV_{(-\infty,\infty)}\Phi(\cdot,t)\leqslant c\omega,\qquad TV_{(-\infty,\infty)}Z(\cdot,t)\leqslant c\omega,\qquad 0\leqslant t<T.
\label{align:ch16-s04-e18}
\end{align}
因此，回忆\eqref{align:ch16-s04-e05}，
\begin{align}
J_s=\sum_{r+s\text{ 奇}}|\hat W_s^r-\hat V_s^r|\leqslant c\rho\omega,\qquad s=0,1,\ldots,s^*.
\label{align:ch16-s04-e19}
\end{align}

下一步是估计 $K_s$，这需要某些准备。当左侧状态 $W$ 与右侧状态 $V$ 由一个可容波扇 $\gamma=(\gamma_1,\ldots,\gamma_n)$ 连接时，
\begin{align}
V=\Omega(\gamma;W),\qquad W=\Theta(\gamma;V),\qquad\gamma=G(W,V),
\label{align:ch16-s04-e20}
\end{align}
其中 $\Omega$ 是由原书(9.3.4)定义的波扇函数，$\Theta(\gamma;\cdot)$ 是固定 $\gamma$ 时 $\Omega(\gamma;\cdot)$ 的反函数，$G(W,\cdot)$ 是固定 $W$ 时 $\Omega(\cdot;W)$ 的反函数（等价地，$G(\cdot,V)$ 是固定 $V$ 时 $\Theta(\cdot;V)$ 的反函数）。
这些函数定义在 $\mathbb{R}^n\times\mathbb{R}^n$ 中原点的某个邻域上，并满足
\begin{align}
\Omega(0;W)=W,\qquad\Theta(0;V)=V,\qquad G(W,W)=0,
\label{align:ch16-s04-e21}\\
\Omega_W(0;W)=I,\qquad\Theta_V(0;V)=I,
\label{align:ch16-s04-e22}\\
\Omega_\gamma(0;W)=R(W),\qquad\Theta_\gamma(0;V)=-R(V),
\label{align:ch16-s04-e23}\\
G_W(W,W)=-L(W),\qquad G_V(V,V)=L(V),
\label{align:ch16-s04-e24}
\end{align}
对定义域中一切 $W$ 与 $V$。代替 $G$，使用函数
\begin{align}
H(W,Y)=G(W,W+Y),
\label{align:ch16-s04-e25}
\end{align}
更为方便，对此我们有
\begin{align}
H(W,0)=0,
\label{align:ch16-s04-e26}\\
H_W(W,0)=0,\qquad H_Y(W,0)=L(W).
\label{align:ch16-s04-e27}
\end{align}

如同第十三章那样，$K_s$ 将由任意固定 $t\in(t_s,t_{s+1})$ 穿越时间线 $t$ 的初等波强度之和来度量。
这些波从网格点 $(x_r,t_s)$（$r+s$ 为偶数）发出。我们将归纳地估计 $K_s$，把它与 $K_{s-1}$ 联系起来。

固定一个网格点 $(x_r,t_s)$，考虑熟悉的菱形区域 $\Delta_s^r$，其顶点为四个周围采样点 $(y_s^{r-1},t_s)$、$(y_{s-1}^r,t_{s-1})$、$(y_s^{r+1},t_s)$ 与 $(y_{s+1}^r,t_{s+1})$，见原书第十三章图 13.3.1。
从 $(x_r,t_s)$ 发出的波扇 $\varepsilon$ 对 $K_s$ 贡献 $|\varepsilon|$。另一方面，从 $(x_{r-1},t_{s-1})$ 发出的波扇的一部分 $\alpha$（可能全部或没有）通过其“西南”边进入 $\Delta_s^r$，从 $(x_{r+1},t_{s-1})$ 发出的波扇的一部分 $\beta$（可能全部或没有）通过其“东南”边进入 $\Delta_s^r$，对 $K_{s-1}$ 贡献 $|\alpha|+|\beta|$。

由 $\hat U_h$ 的构造，结合\eqref{align:ch16-s04-e20}与\eqref{align:ch16-s04-e25}，
\begin{align}
U_s^{r-1}=\Theta(\alpha;V_{r-1}^s)+\Phi_s^{r-1}-\Phi_{s-1}^{r-1},
\label{align:ch16-s04-e28}\\
U_s^{r+1}=\Omega(\beta;W_{s-1}^r)+\Phi_s^{r+1}-\Phi_{s-1}^{r+1},
\label{align:ch16-s04-e29}\\
\varepsilon=G(W_s^{r-1},V_s^{r+1})=H(W_s^{r-1},V_s^{r+1}-W_s^{r-1}).
\label{align:ch16-s04-e30}
\end{align}
我们现在定义
\begin{align}
\tilde U_s^{r-1}=\Theta(\alpha;U_{s-1}^r),
\label{align:ch16-s04-e31}\\
\tilde U_s^{r+1}=\Omega(\beta;U_{s-1}^r),
\label{align:ch16-s04-e32}\\
\tilde\varepsilon=G(\tilde U_s^{r-1},\tilde U_s^{r+1})=H(\tilde U_s^{r-1},\tilde U_s^{r+1}-\tilde U_s^{r-1}).
\label{align:ch16-s04-e33}
\end{align}
由原书(13.3.1)，
\begin{align}
|\tilde\varepsilon-\alpha+\beta|\leqslant c\mathscr{D}(\Delta_s^r),
\label{align:ch16-s04-e34}
\end{align}
其中 $\mathscr{D}(\Delta_s^r)$ 由原书(13.3.2)给出。

下一步是估计 $\varepsilon-\tilde\varepsilon$。由\eqref{align:ch16-s04-e30}与\eqref{align:ch16-s04-e33}，
\begin{align}
\varepsilon-\tilde\varepsilon=\bar H_W[W_s^{r-1}-\tilde U_s^{r-1}]+\bar H_Y[U_s^{r+1}-W_s^{r-1}-\tilde U_s^{r+1}+\tilde U_s^{r-1}]
\label{align:ch16-s04-e35}
\end{align}
其中
\begin{align}
\bar H_W=\int_0^1 H_W(\xi W_s^{r-1}+(1-\xi)\tilde U_m^{r-1},\tilde U_s^{r+1}-\tilde U_s^{r-1})\,\mathrm{d}\xi,
\label{align:ch16-s04-e36}\\
\bar H_Y=\int_0^1 H_Y(W_s^{r-1},\xi(V_s^{r+1}-W_s^{r-1})+(1-\xi)(\tilde U_s^{r+1}-\tilde U_s^{r-1}))\,\mathrm{d}\xi.
\label{align:ch16-s04-e37}
\end{align}
合并\eqref{align:ch16-s04-e07}、\eqref{align:ch16-s04-e28}、\eqref{align:ch16-s04-e29}、\eqref{align:ch16-s04-e31}与\eqref{align:ch16-s04-e32}，
\begin{align}
W_s^{r-1}-\tilde U_s^{r-1}&=\bar\Theta_V\{Q_{s-1}^r[\Phi_{s-1}^{r-1}-\Phi_{s-1}^r]-\tau P_{s-1}^r\}\nonumber\\
&\quad+\Phi_s^{r-1}-\Phi_{s-1}^{r-1}+Q_s^{r-1}[\Phi_s^r-\Phi_{s-1}^{r-1}]-\tau P_s^{r-1},
\label{align:ch16-s04-e38}\\
V_s^{r+1}-\tilde U_s^{r+1}&=\bar\Omega_W\{Q_{s-1}^r[\Phi_{s-1}^{r+1}-\Phi_{s-1}^r]-\tau P_{s-1}^r\}\nonumber\\
&\quad+\Phi_s^{r+1}-\Phi_{s-1}^{r+1}+Q_s^{r+1}[\Phi_s^r-\Phi_s^{r+1}]-\tau P_s^{r+1},
\label{align:ch16-s04-e39}
\end{align}
其中
\begin{align}
\bar\Theta_V=\int_0^1\Theta_V(\alpha;\xi V_{s-1}^r+(1-\xi)U_{s-1}^r)\,\mathrm{d}\xi,
\label{align:ch16-s04-e40}\\
\bar\Omega_W=\int_0^1\Omega_W(\beta;\xi W_{s-1}^r+(1-\xi)U_{s-1}^r)\,\mathrm{d}\xi.
\label{align:ch16-s04-e41}
\end{align}

由原书(14.6.36)、\eqref{align:ch16-s04-e27}、\eqref{align:ch16-s04-e31}、\eqref{align:ch16-s04-e21}与\eqref{align:ch16-s04-e32}，我们推得 $|\bar H_W|\leqslant c(|\alpha|+|\beta|)$。
此外，\eqref{align:ch16-s04-e38}、\eqref{align:ch16-s04-e03}、\eqref{align:ch16-s03-e08}、\eqref{align:ch16-s04-e06}与\eqref{align:ch16-s03-e12}给出 $|W_s^{r-1}-\tilde U_s^{r-1}|\leqslant c\rho h$。因此，我们可以如下界住\eqref{align:ch16-s04-e35}右端的第一项：
\begin{align}
|\bar H_W[W_s^{r-1}-\tilde U_s^{r-1}]|\leqslant c\rho(|\alpha|+|\beta|)h.
\label{align:ch16-s04-e42}
\end{align}

\eqref{align:ch16-s04-e35}右端的第二项需要更精细的处理。由\eqref{align:ch16-s04-e38}与\eqref{align:ch16-s04-e39}，我们推得
\begin{align}
&W_s^{r-1}-\tilde U_s^{r-1}-V_s^{r+1}+\tilde U_s^{r+1}\nonumber\\
&=\tau[P_s^{r+1}-P_s^{r-1}]\nonumber\\
&\quad+[Q_{s-1}^{r-1}-I][\Phi_s^{r+1}-\Phi_{s-1}^{r+1}-\Phi_s^{r-1}+\Phi_{s-1}^{r-1}]\nonumber\\
&\quad+[\bar\Theta_V-I]\{Q_{m-1}^r[\Phi_{s-1}^{r-1}-\Phi_{s-1}^r]-\tau P_{s-1}^r\}\nonumber\\
&\quad-[\bar\Omega_W-I]\{Q_{m-1}^r[\Phi_{s-1}^{r+1}-\Phi_{s-1}^r]-\tau P_{s-1}^r\}\nonumber\\
&\quad+[Q_s^{r-1}-Q_{s-1}^r][\Phi_s^r-\Phi_{s-1}^{r-1}]\nonumber\\
&\quad-[Q_s^{r+1}-Q_{s-1}^r][\Phi_s^r-\Phi_s^{r+1}].
\label{align:ch16-s04-e43}
\end{align}
回忆\eqref{align:ch16-s04-e06}与\eqref{align:ch16-s03-e12}，
\begin{align}
P_s^{r+1}-P_s^{r-1}=\bar{\hat P}_{\hat U}[\hat U_s^{r+1}-\hat U_s^{r-1}]+\bar{\hat P}_\Phi[\Phi_s^{r+1}-\Phi_s^{r-1}]-[Z_s^{r+1}-Z_s^{r-1}],
\label{align:ch16-s04-e44}
\end{align}
其中
\begin{align}
\bar{\hat P}_{\hat U}=\int_0^1\hat P_{\hat U}(\xi\hat U_s^{r+1}+(1-\xi)\hat U_s^{r-1},\Phi_s^{r+1},0)\,\mathrm{d}\xi,
\label{align:ch16-s04-e45}\\
\bar{\hat P}_\Phi=\int_0^1\hat P_\Phi(\hat U_s^{r-1},\xi\Phi_s^{r+1}+(1-\xi)\Phi_s^{r-1},0)\,\mathrm{d}\xi.
\label{align:ch16-s04-e46}
\end{align}
此外，合并\eqref{align:ch16-s04-e04}、\eqref{align:ch16-s04-e28}、\eqref{align:ch16-s04-e29}、\eqref{align:ch16-s04-e21}、\eqref{align:ch16-s04-e07}并作简短计算，我们得到
\begin{align}
\hat U_s^{r+1}-\hat U_s^{r-1}=\bar\Omega_\gamma\beta-\bar\Theta_\gamma\alpha+[Q_{s-1}^r-I][\Phi_{s-1}^{r+1}-\Phi_{s-1}^{r-1}],
\label{align:ch16-s04-e47}
\end{align}
其中
\begin{align}
\bar\Omega_\gamma=\int_0^1\Omega_\gamma(\xi\beta;W_{s-1}^r)\,\mathrm{d}\xi,\qquad\bar\Theta_\gamma=\int_0^1\Theta_\gamma(\xi\alpha;V_{s-1}^{r})\,\mathrm{d}\xi.
\label{align:ch16-s04-e48}
\end{align}

继续估计\eqref{align:ch16-s04-e43}右端的剩余项，我们使用\eqref{align:ch16-s04-e22}、\eqref{align:ch16-s04-e40}与\eqref{align:ch16-s04-e41}得到
\begin{align}
|\bar\Theta_V-I|\leqslant c|\alpha|,\qquad|\bar\Omega_W-I|\leqslant c|\beta|.
\label{align:ch16-s04-e49}
\end{align}
此外，使用\eqref{align:ch16-s04-e05}、\eqref{align:ch16-s04-e03}、\eqref{align:ch16-s04-e28}、\eqref{align:ch16-s04-e07}、\eqref{align:ch16-s04-e06}、\eqref{align:ch16-s03-e12}与\eqref{align:ch16-s03-e08}，
\begin{align}
|Q_s^{r-1}-Q_{s-1}^r|\leqslant c|\alpha|+c\rho h,\qquad|Q_s^{r+1}-Q_{s-1}^r|\leqslant c|\beta|+c\rho h.
\label{align:ch16-s04-e50}
\end{align}

现在合并\eqref{align:ch16-s04-e34}、\eqref{align:ch16-s04-e35}、\eqref{align:ch16-s04-e42}、\eqref{align:ch16-s04-e43}、\eqref{align:ch16-s04-e44}、\eqref{align:ch16-s04-e47}、\eqref{align:ch16-s04-e49}与\eqref{align:ch16-s04-e50}，推出
\begin{align}
|\varepsilon-\alpha-\beta|\leqslant c\rho(|\alpha|+|\beta|)h+c\mathscr{D}(\Delta_s^r)\nonumber\\
\qquad+c\{|\Phi_s^{r+1}-\Phi_s^r|+|\Phi_s^r-\Phi_s^{r-1}|+|Z_s^{r+1}-Z_s^{r-1}|\}h\nonumber\\
\qquad+c\rho|\Phi_s^{r+1}-\Phi_{s-1}^{r+1}-\Phi_s^{r-1}+\Phi_{s-1}^{r-1}|,
\label{align:ch16-s04-e51}\\
|\varepsilon|\leqslant|I+\tau\bar H_Y\bar{\hat P}_{\hat U}\bar\Theta_\gamma||\alpha|+|I-\tau\bar H_Y\bar{\hat P}_{\hat U}\bar\Omega_\gamma||\beta|\nonumber\\
\qquad+c\rho(|\alpha|+|\beta|)h+c\mathscr{D}(\Delta_s^r)\nonumber\\
\qquad+c\{|\Phi_s^{r+1}-\Phi_s^r|+|\Phi_s^r-\Phi_s^{r-1}|+|Z_s^{r+1}-Z_s^{r-1}|\}h\nonumber\\
\qquad+c\rho|\Phi_s^{r+1}-\Phi_{s-1}^{r+1}-\Phi_s^{r-1}+\Phi_{s-1}^{r-1}|.
\label{align:ch16-s04-e52}
\end{align}

当 $\rho$ 充分小时，由\eqref{align:ch16-s04-e37}、\eqref{align:ch16-s04-e27}、\eqref{align:ch16-s04-e48}、\eqref{align:ch16-s04-e23}与\eqref{align:ch16-s04-e45}，$\bar H_Y$ 接近 $L(0)$，$\bar\Omega_\gamma$ 接近 $R(0)$，$\bar\Theta_\gamma$ 接近 $-R(0)$，$\bar{\hat P}_{\hat U}$ 接近 $\hat P_{\hat U}(0,0,0)$。因此，$-\bar H_Y\bar{\hat P}_{\hat U}\bar\Theta_\gamma$ 与 $\bar H_Y\bar{\hat P}_{\hat U}\bar\Omega_\gamma$ 二者都接近由\eqref{align:ch16-s03-e13}定义的矩阵 $\hat A$。特别地，这些矩阵列对角占优，从而
\begin{align}
|I+\tau\bar H_Y\bar{\hat P}_{\hat U}\bar\Theta_\gamma|\leqslant1-4\nu\tau,\qquad|I-\tau\bar H_Y\bar{\hat P}_{\hat U}\bar\Omega_\gamma|\leqslant1-4\nu\tau,
\label{align:ch16-s04-e53}
\end{align}
对某个 $\nu>0$。于是，当 $\rho$ 充分小时，\eqref{align:ch16-s04-e52}给出估计
\begin{align}
|\varepsilon|\leqslant(1-3\nu\tau)(|\alpha|+|\beta|)+c\mathscr{D}(\Delta_s^r)\nonumber\\
\qquad+c\{|\Phi_s^{r+1}-\Phi_s^r|+|\Phi_s^r-\Phi_s^{r-1}|+|Z_s^{r+1}-Z_s^{r-1}|\}h\nonumber\\
\qquad+c\rho|\Phi_s^{r+1}-\Phi_{s-1}^{r+1}-\Phi_s^{r-1}+\Phi_{s-1}^{r-1}|.
\label{align:ch16-s04-e54}
\end{align}

固定 $s$，对所有 $r+s$ 为偶数者对不等式\eqref{align:ch16-s04-e54}求和。回顾\eqref{align:ch16-s04-e16}，我们注意到 $|\varepsilon|$ 项之和给出 $K_s$，而 $|\alpha|+|\beta|$ 项之和给出 $K_{s-1}$。$|\Phi_s^{r+1}-\Phi_s^r|$、$|\Phi_s^r-\Phi_s^{r-1}|$ 与 $|Z_s^{r+1}-Z_s^{r-1}|$ 的和由 $c\omega$ 界住，这由\eqref{align:ch16-s04-e03}与\eqref{align:ch16-s04-e18}得到。最后，合并\eqref{align:ch16-s04-e03}与\eqref{align:ch16-s04-e08}，
\begin{align}
&\Phi_s^{r+1}-\Phi_{s-1}^{r+1}-\Phi_s^{r-1}+\Phi_{s-1}^{r-1}\nonumber\\
&=-\int_{t_{s-1}}^{t_s}N[F(U(x_{r+1},t))-F(U(x_{r-1},t))]\,\mathrm{d}t-\int_{t_{s-1}}^{t_s}[Z(x_{r+1},t)-Z(x_{r-1},t)]\,\mathrm{d}t,
\label{align:ch16-s04-e55}
\end{align}
由此，由\eqref{align:ch16-s04-e01}与\eqref{align:ch16-s04-e18}，
\begin{align}
\sum_{r+s\text{ 偶}}|\Phi_s^{r+1}-\Phi_{s-1}^{r+1}-\Phi_s^{r-1}+\Phi_{s-1}^{r-1}|\leqslant c(\omega+X_s)\tau,
\label{align:ch16-s04-e56}
\end{align}
其中
\begin{align}
X_s=\frac{1}{\tau}\int_{t_{s-1}}^{t_s}X(t)\,\mathrm{d}t.
\label{align:ch16-s04-e57}
\end{align}
于是我们有
\begin{align}
K_s\leqslant(1-3\nu\tau)K_{s-1}+c\omega\tau+c\rho X_s\tau+c\Xi_s,
\label{align:ch16-s04-e58}
\end{align}
其中
\begin{align}
\Xi_s=\sum_{r+s\text{ 偶}}\mathscr{D}(\Delta_s^r).
\label{align:ch16-s04-e59}
\end{align}

为估计项 $\Xi_s$，我们采用第十三章第 4 节中使用的程序。
我们引入 Glimm 泛函
\begin{align}
\Gamma_s=K_s+\kappa M_s,
\label{align:ch16-s04-e60}
\end{align}
其中 $\kappa$ 是一个待在下面固定的正数，而
\begin{align}
M_s=\sum|\gamma_i||\theta_j|,
\label{align:ch16-s04-e61}
\end{align}
求和遍历从网格点 $(x_r,t_s)$（$r+s$ 为偶数）发出的所有相互靠近的初等波对 $(\gamma_i,\theta_j)$。

显然，$K_s\leqslant\Gamma_s\leqslant(1+c\rho)K_s$。此外，
\begin{align}
M_s\leqslant M_{s-1}+(K_s+K_{s-1})E_s-\Xi_s,
\label{align:ch16-s04-e62}
\end{align}
其中 $E_s$ 表示与 $r+s$ 为偶数的 $\Delta_s^r$ 相关联的 $|\varepsilon-\alpha-\beta|$ 之和。在对所有 $r+s$ 为偶数者对不等式\eqref{align:ch16-s04-e51}求和之后，我们推得
\begin{align}
E_s\leqslant c\rho K_{s-1}\tau+c\omega\tau+c\rho X_s.
\label{align:ch16-s04-e63}
\end{align}
因此，鉴于\eqref{align:ch16-s04-e60}、\eqref{align:ch16-s04-e58}、\eqref{align:ch16-s04-e62}与\eqref{align:ch16-s04-e63}，我们推出：当 $\kappa$ 大且 $\rho$ 小时，
\begin{align}
\Gamma_s\leqslant(1-2\nu\tau)\Gamma_{s-1}+c\omega\tau+c\rho X_s\tau+c\Xi_s.
\label{align:ch16-s04-e64}
\end{align}
对\eqref{align:ch16-s04-e64}迭代 $s=1,2,\ldots$，我们得到
\begin{align}
\Gamma_s\leqslant(1-2\nu\tau)^s\Gamma_0+c\omega+c\rho\tau\sum_{k=1}^s(1-2\nu\tau)^{s-k}X_k.
\label{align:ch16-s04-e65}
\end{align}
回忆\eqref{align:ch16-s04-e15}、\eqref{align:ch16-s04-e16}、\eqref{align:ch16-s04-e19}与\eqref{align:ch16-s03-e15}，我们推出
\begin{align}
\hat X_s\leqslant c(1-2\nu\tau)^s\delta+c\omega+c\rho\tau\sum_{k=1}^s(1-2\nu\tau)^{s-k}X_k.
\label{align:ch16-s04-e66}
\end{align}
于是，对任意 $t\in[0,T)$，$\hat U_h(\cdot,t)$ 的总变差关于 $h$ 一致有界。

如同第十三章那样，为建立约似解族 $\{\hat U_h\}$ 的紧性，我们还需要 $t$ 方向的等度连续性。
为此，固定 $\ell>1$。对任意 $\zeta,\xi\in[t_s,t_{s+1}]$，$\zeta<\xi$：
\begin{align}
\int_{-\ell}^{\ell}|\hat U_h(x,\xi)-\hat U_h(x,\zeta)|\,\mathrm{d}x\leqslant c\hat X_s(\xi-\zeta)\leqslant c\rho\tau.
\label{align:ch16-s04-e67}
\end{align}
回忆\eqref{align:ch16-s04-e10}、\eqref{align:ch16-s04-e09}与\eqref{align:ch16-s04-e12}，
\begin{align}
&\int_{-\ell}^{\ell}|\hat U_h(x,t_s)-\hat U_h(x,t_s-)|\,\mathrm{d}x\nonumber\\
&\qquad\leqslant\sum\int_{x_{r-1}}^{x_{r+1}}|\hat U_h(y_s^r,t_s-)-\hat U_h(x,t_s-)|\,\mathrm{d}x\nonumber\\
&\qquad\quad+\sum h|Q_s^r-I||\Phi_s^{r+1}+\Phi_s^{r-1}-2\Phi_s^r|+\sum 2h\tau|\hat P_s^r|+O(h),
\label{align:ch16-s04-e68}
\end{align}
其中求和遍历所有满足 $(|r|+1)h<\ell$ 且 $r+s$ 为奇数的 $r$。在\eqref{align:ch16-s04-e68}的右端，第一个和由 $2h\hat X_{s-1}$ 界住，第二个和被 $c\rho h\omega$ 控制，最后一个和由 $c\rho\ell\tau$ 界住。因此，合并\eqref{align:ch16-s04-e67}与\eqref{align:ch16-s04-e68}，我们推得：对任意 $0\leqslant\zeta<\xi<T$，
\begin{align}
\int_{-\ell}^{\ell}|\hat U_h(x,\xi)-\hat U_h(x,\zeta)|\,\mathrm{d}x\leqslant c\ell(|\xi-\zeta|+h).
\label{align:ch16-s04-e69}
\end{align}

于是我们建立了族 $\{\hat U_h\}$ 所需的紧性，可以抽取子列 $\{\hat U_{h_m}\}$（$h_m\to0$，当 $m\to\infty$），它对每个 $t\in[0,\infty)$ 和几乎所有 $x\in(-\infty,\infty)$ 收敛到一个 $BV$ 函数 $\hat U$。我们现在证明 $\hat U$ 是以\eqref{align:ch16-s03-e15}给出的 $\hat U_0$ 为初值的\eqref{align:ch16-s03-e10}的解。

我们固定任意光滑试验函数 $\psi$，在 $(-\infty,\infty)\times[0,T)$ 上具紧支集。以 $\psi$ 乘\eqref{align:ch16-s04-e13}，在由\eqref{align:ch16-s04-e09}标识的矩形 $\mathscr{R}_s^r$ 上积分，分部积分，并对 $r$ 与 $s$（$r+s$ 为偶数）求和。这给出
\begin{align}
&\sum_{s=0}^{s^*}\sum_{r+s\text{ 偶}}\int_{t_s}^{t_{s+1}}\int_{x_{r-1}}^{x_{r+1}}[\partial_t\psi\hat U_h+\partial_x\psi\hat F(\hat U_h,\Phi_s^r)]\,\mathrm{d}x\,\mathrm{d}t\nonumber\\
&\quad+\sum_{s=0}^{s^*}\sum_{r+s\text{ 奇}}\Bigl[\int_{t_s}^{t_{s+1}}\psi(x_r,t)\,\mathrm{d}t\Bigr][\hat F(\hat W_s^r,\Phi_s^{r+1})-\hat F(\hat V_s^r,\Phi_s^{r-1})]\nonumber\\
&\quad+\sum_{s=0}^{s^*}\int_{-\infty}^{\infty}\psi(x,t_s)[\hat U_h(x,t_s)-\hat U_h(x,t_s-)]\,\mathrm{d}x+\int_{-\infty}^{\infty}\psi(x,0)\hat U_0(x)\,\mathrm{d}x=0.
\label{align:ch16-s04-e070}
\end{align}

我们考察上式左端每一项当沿序列 $\{h_m\}$ 令 $h\to0$ 时的行为。第一项收敛到
\begin{align}
\int_0^T\int_{-\infty}^{\infty}[\partial_t\psi\hat U+\partial_x\psi\hat F(\hat U,\Phi)]\,\mathrm{d}x\,\mathrm{d}t.
\label{align:ch16-s04-e071}
\end{align}
第二项是 $O(h)$，因为由\eqref{align:ch16-s04-e07}与\eqref{align:ch16-s04-e08}，
\begin{align}
\hat F(\hat W_s^r,\Phi_s^{r+1})-\hat F(\hat V_s^r,\Phi_s^{r-1})=O(h^2).
\label{align:ch16-s04-e072}
\end{align}
对第三项，我们使用\eqref{align:ch16-s04-e14}、\eqref{align:ch16-s04-e08}、\eqref{align:ch16-s04-e07}与\eqref{align:ch16-s04-e04}得到
\begin{align}
&\int_{-\infty}^{\infty}\psi(x,t_s)[\hat U_h(x,t_s)-\hat U_h(x,t_s-)]\,\mathrm{d}x\nonumber\\
&=\sum_{r+s\text{ 奇}}\int_{x_{r-1}}^{x_{r+1}}\psi(x,t_s)[\hat U_h(y_s^r,t_s-)-\hat U_h(x,t_s-)]\,\mathrm{d}x\nonumber\\
&\quad+\sum_{r+s\text{ 奇}}h\psi(x_r,t_s)[Q_s^r-I][\Phi_s^{r+1}+\Phi_s^{r-1}-2\Phi_s^r]\nonumber\\
&\quad+\sum_{r+s\text{ 奇}}\Bigl\{\int_{x_{r-1}}^{x_r}[\psi(x,t_s)-\psi(x_r,t_s)]\,\mathrm{d}x\Bigr\}[Q_s^r-I][\Phi_s^{r-1}-\Phi_s^r]\nonumber\\
&\quad+\sum_{r+s\text{ 奇}}\Bigl\{\int_{x_r}^{x_{r+1}}[\psi(x,t_s)-\psi(x_r,t_s)]\,\mathrm{d}x\Bigr\}[Q_s^r-I][\Phi_s^{r+1}-\Phi_s^r]\nonumber\\
&\quad-\sum_{r+s\text{ 奇}}\tau\Bigl\{\int_{x_{r-1}}^{x_{r+1}}\psi(x,t_s)\,\mathrm{d}x\Bigr\}\hat P_s^r.
\label{align:ch16-s04-e073}
\end{align}

我们接着估计上式右端每一项，着眼于它们对\eqref{align:ch16-s04-e070}的贡献。第二项由 $c\rho h^2X(s)$ 界住，而第三与第四项都由 $c\rho h^2\omega$ 界住。于是上述三项对\eqref{align:ch16-s04-e070}的贡献是 $O(h)$。相反，\eqref{align:ch16-s04-e073}右端的第一项由 $ch\hat X_{s-1}$ 界住，因而表面上可能对\eqref{align:ch16-s04-e070}贡献 $O(1)$。然而，几乎可以肯定，该项的贡献是 $o(1)$，理由正如第十三章中所讨论的：采样点 $(y_s^r,t_s)$ 是随机选取的。
至于\eqref{align:ch16-s04-e073}中的最后一项，我们有
\begin{align}
\left|\tau\int_{x_{r-1}}^{x_{r+1}}\psi(x,t_s)\hat P_s^r\,\mathrm{d}x-\int_{t_{s-1}}^{t_s}\int_{x_{r-1}}^{x_{r+1}}\psi(x,t)\hat P(\hat U_h(x,t),\Phi(x,t),Z(x,t))\,\mathrm{d}x\,\mathrm{d}t\right|\leqslant ch^2(u_{s-1}^r+h),
\label{align:ch16-s04-e074}
\end{align}
其中 $u_{s-1}^r$ 表示 $\hat U_h$ 在矩形 $\mathscr{R}_{s-1}^r$ 上的振荡。注意，由 $\hat U_h$ 的构造，它在 $\mathscr{R}_{s-1}^r$ 上的振荡由 $\hat U_h(\cdot,t_s-)$ 在区间 $(x_{r-1},x_{r+1})$ 上的变差界住，从而所有 $r+s$ 为奇数者上 $u_{s-1}^r$ 之和由 $\hat X_{s-1}$ 界住。于是\eqref{align:ch16-s04-e073}中最后一项对\eqref{align:ch16-s04-e070}的贡献是
\begin{align}
-\int_0^T\int_{-\infty}^{\infty}\psi\hat P(\hat U_h,\Phi,Z)\,\mathrm{d}x\,\mathrm{d}t+O(h).
\label{align:ch16-s04-e075}
\end{align}

现在我们合并以上全部，并在\eqref{align:ch16-s04-e070}中取 $h\to0$ 极限，得到
\begin{align}
\int_0^T\int_{-\infty}^{\infty}[\partial_t\psi\hat U+\partial_x\psi\hat F(\hat U,\Phi)-\psi\hat P(\hat U,\Phi,Z)]\,\mathrm{d}x\,\mathrm{d}t\nonumber\\
\qquad+\int_{-\infty}^{\infty}\psi(x,0)\hat U_0(x)\,\mathrm{d}x=0,
\label{align:ch16-s04-e076}
\end{align}
这验证了 $\hat U$ 是以 $\hat U_0$ 为初值的\eqref{align:ch16-s03-e10}的解。

在令
\begin{align}
\hat X(t)=TV_{(-\infty,\infty)}\hat U(\cdot,t),\qquad 0\leqslant t<T
\label{align:ch16-s04-e077}
\end{align}
之后，\eqref{align:ch16-s04-e66}给出估计
\begin{align}
\hat X(t)\leqslant c\delta e^{-2\nu t}+c\omega+c\rho\int_0^t e^{-2\nu(t-\xi)}X(\xi)\,\mathrm{d}\xi.
\label{align:ch16-s04-e078}
\end{align}

在下文，我们将默认\eqref{align:ch16-s03-e10}的 Cauchy 问题拥有唯一可容 $BV$ 解。证明这一点需要把第 14.10 节中从齐次守恒律组到非齐次平衡律组的分析推广——这是一项不应涉及主要新思想的艰巨任务。
在唯一性的预设下，由于 $U-\Phi$ 是以\eqref{align:ch16-s03-e15}定义的 $\hat U_0$ 为初值的\eqref{align:ch16-s03-e10}的解，我们断定 $\hat U=U-\Phi$，特别地 $X(t)\leqslant\hat X(t)+c\omega$。因此，\eqref{align:ch16-s04-e078}给出
\begin{align}
X(t)\leqslant c\delta e^{-2\nu t}+c\omega+c\rho\int_0^t e^{-2\nu(t-\xi)}X(\xi)\,\mathrm{d}\xi.
\label{align:ch16-s04-e079}
\end{align}
我们令
\begin{align}
Y(t)=\int_0^t e^{2\nu\xi}X(\xi)\,\mathrm{d}\xi,
\label{align:ch16-s04-e080}
\end{align}
此时\eqref{align:ch16-s04-e079}给出微分不等式
\begin{align}
\dot Y(t)\leqslant c\rho Y(t)+c\delta+c\omega e^{2\nu t}.
\label{align:ch16-s04-e081}
\end{align}
我们取 $\rho$ 充分小，使得乘 $Y(t)$ 的系数 $c\rho$ 小于 $\nu$。然后积分\eqref{align:ch16-s04-e081}并代回\eqref{align:ch16-s04-e079}，得到
\begin{align}
X(t)\leqslant c\omega+c\delta e^{-\nu t},\qquad 0\leqslant t\leqslant T,
\label{align:ch16-s04-e082}
\end{align}
这建立了估计\eqref{align:ch16-s03-e03}，并完成了\cref{theorem:ch16-1631}的证明。

\end{document}
